\documentclass[11pt,a4paper]{article}

\usepackage[margin=2.4cm]{geometry}
\usepackage[T1]{fontenc}
\usepackage[utf8]{inputenc}
\usepackage{microtype}
\usepackage{amsmath,amssymb}
\usepackage{booktabs}
\usepackage{xcolor}
\usepackage{caption}
\usepackage{graphicx}
\usepackage{pgfplots}
\usepackage[numbers,sort&compress]{natbib}
\usepackage[hidelinks]{hyperref}
\usepackage{orcidlink}

\pgfplotsset{compat=1.18}
\definecolor{cReal}{HTML}{1B4965}
\definecolor{cBeta}{HTML}{C1666B}
\definecolor{cShock}{HTML}{7C9885}
\definecolor{cOne}{HTML}{D08C34}

\captionsetup{font=small,labelfont=bf}
\newcommand{\neff}{n_{\mathrm{eff}}}
\newcommand{\code}[1]{\texttt{#1}}

\title{\bfseries Reviewer Count Is Not Evidence Count:\\
Measured Error Dependence in an Executable Verification Panel}

\author{Mohsen Seyedkazemi Ardebili\\
  \small\href{https://mskazemi.github.io}{mskazemi.github.io}\\
  \small\texttt{mohsen.seyedkazemi@gmail.com}}

\date{\today}

\begin{document}
\maketitle

\begin{abstract}
\noindent
Panels of automated verifiers --- test oracles, static analysers, and increasingly
LLM reviewers --- are routinely sized by head count, on the implicit assumption
that adding a verifier adds evidence. We measure that assumption on a panel whose
ground truth is decided by execution rather than by judgement: 25 partial test
oracles voting on 72 program candidates drawn from 24 problems, where a
candidate's viability is the exit status of a hidden test suite. Every verifier
discriminates significantly above chance (mean accuracy $p = 0.7956$, 25/25
significant Youden's $J$ after Bonferroni correction), so the panel's dependence
structure is interpretable rather than an artifact of incompetent members.

Three results follow. First, the panel buys nothing over its own members:
majority-vote error is $0.2083$ (95\% CI $[0.131, 0.316]$) against a mean
member's $0.2044$ and a \emph{best} member's $0.1389$ --- so the panel is not
worth one verifier, it is worse than one, and its effective independent size
falls below the metric's floor of 1 rather than landing on it. On 72 items that
ordering is directional, not significant (panel vs.\ best member, exact McNemar
$p = 0.18$), and we report it as such. The design-effect heuristic
$N/(1+(N-1)\rho)$, widely reused to discount correlated voters, reports $1.66$
for this panel; under the dependence shape our data support it overstates
available independence in $441/441$ cells of a $63{,}000$-run simulation grid,
though a contemporaneous LLM-judge panel finds the same heuristic accurate at a
lower correlation, so the failure is regime-specific rather than universal. Second, that shape is not the shared-shock
mixture that recent multi-agent work assumes. Real panels fail \emph{partially}
--- 11 of the 15 observed panel failures had a majority wrong while a
minority was still right --- an outcome shared shock assigns essentially zero
probability. A per-item difficulty model in the family proposed by
Eckhardt and Lee in 1985, at identical parameter count, reproduces the partial
failures ($11.1$ predicted against $11$ observed) and lands $3.7\times$ closer on
panel error. Third, neither two-parameter model gets the high-quorum limit right,
because the real panel has an \emph{irreducible blind spot}: four of 72
candidates are missed by all 25 verifiers, a floor of $\lambda = 0.0556$
(95\% CI $[0.022, 0.134]$) that no acceptance threshold and no panel size can
cross. Shared shock over-predicts
that floor by $2.13\times$ and per-item difficulty promises it away entirely.

The practical consequence is that a verification panel must be characterised by
measurement, not by declared diversity: verifiers sharing no declared attribute
in our panel still shared $53\%$ of their errors, and weighting votes by
declared independence groups --- an intuitively safe defence --- never wins when
that declared structure is not real. Every finding rests on one constructed
panel of 25 oracles over 72 items; all code, artifacts, and the analysis
pipeline are public.
\end{abstract}

\section{Introduction}

Software increasingly arrives faster than people can review it. The response,
across code review, CI, benchmark scoring, and agentic coding pipelines, has been
to add reviewers: more static analysers, more test oracles, more model instances
voting on the same artifact, aggregated by majority or by quorum. The reasoning
is the Condorcet jury theorem's: individually fallible voters, combined, are
reliably right.

That theorem's premise is independence, which both its classical statement and
its generalisations call unrealistic --- Condorcet's assumptions of independence
and homogeneity are ``seldom realistic'' \citep{boland1989}, and the theorem
``is known to hold only under the unrealistic assumption that votes are
independent'' \citep{ladha1992}. Both papers relax that premise rather than
abandon it, deriving conditions under which majority voting survives dependence. Software engineering learned this the hard
way once already. N-version programming \citep{avizienis1977} runs multiple
independently developed implementations of one specification and votes on their
outputs; in the words of the experiment that tested it, the method ``depends for
its reliability improvement on the assumption that programs that have been
developed independently will fail independently'' \citep{knight1986}.
\citet{knight1986} then tested exactly that with 27 versions written by different
people from the same specification, and found the coincident failure rate far
above what independence predicts. The theoretical accounts that followed
\citep{eckhardt1985,littlewood1989} replaced the independence assumption with a
per-input difficulty function: some inputs are simply hard, and every version
tends to fail on them together.

The panels being assembled today are not independently written programs by
different teams, but they inherit the same question and answer it less carefully.
A modern verification panel is often five prompts against one model family, or
five analysers sharing a parser, or five test oracles sampling one input domain.
Its head count is reported; its independence is assumed. Where a discount is
applied at all, it is the survey-sampling design effect
$\neff = N/(1 + (N-1)\rho)$ \citep{kish1965}, carried over from cluster sampling
and evaluated at some estimate of inter-rater agreement ---
\citet{kohli2026} apply it to an LLM-judge panel in exactly this way.

This paper asks the quantitative version of the question on a panel where the
answer can be checked:

\begin{quote}
Given $n$ verifiers with per-verifier accuracy $p$ and measured error
correlation $\rho$, how many statistically independent verifiers is that panel
actually worth --- and does any standard model of verifier dependence predict
what a real panel does?
\end{quote}

\paragraph{Contributions.}
\begin{enumerate}
\item \textbf{An instrument where dependence is measurable} (\S\ref{sec:instrument}).
  A 25-verifier panel of partial test oracles over 72 candidates with
  execution-decided ground truth. We first report the negative result that
  motivated it: a 20-member panel of four small open-weight LLMs across five
  prompt templates produced no member that discriminated above chance, making
  its correlation structure uninterpretable.
\item \textbf{A panel that loses to its own members} (\S\ref{sec:measure}).
  Majority vote errs on $0.2083$ of items against the best member's $0.1389$, so
  the panel's effective independent size falls \emph{below} the metric's floor of
  1. Declared independence labels carry real but partial signal ($\rho = 0.892$
  within group against $0.526$ across). \citet{kohli2026} reported the same
  headline on an LLM-judge panel while this work was in preparation, and has
  priority on that framing; \S\ref{sec:related-panels} reconciles the two,
  including where they disagree.
\item \textbf{A model-shape result} (\S\ref{sec:shape}). The shared-shock mixture
  is not merely mis-parameterised but structurally wrong: it cannot produce
  partial panel failure, which is the dominant real failure mode. A per-item
  difficulty (beta-binomial) model at equal parameter count fits.
\item \textbf{An irreducible blind spot} (\S\ref{sec:quorum}). The real quorum
  frontier bottoms out at $\lambda = 0.0556$ and stops. Neither two-parameter
  model expresses this; a one-inflated extension does. Escaping a correlation and
  escaping a blind spot require different interventions.
\item \textbf{Design consequences} (\S\ref{sec:consequences}): a negative result
  on group-balanced aggregation, and a single-factor decomposition showing that
  panel \emph{correlation} --- not accuracy, and not the blind spot --- carries
  $53$ of the $58$ catastrophic outcomes an optimising adversary induces.
\end{enumerate}

Everything reported here is reproducible from committed artifacts in the IDKMesh
repository \citep{idkmesh}; \S\ref{sec:repro} gives the commands.

\section{Background and related work}
\label{sec:related}

\paragraph{Correlated failure in redundant software.}
The closest prior work is the multiversion programming literature.
\citet{knight1986} is the empirical anchor: independence, assumed as a design
property, did not hold. \citet{eckhardt1985} give the model we find fits ---
inputs carry a difficulty $\theta(x)$, versions fail independently
\emph{conditional} on $x$, and the resulting marginal failures are positively
correlated. \citet{littlewood1989} extend it to versions built by deliberately
different methodologies.

Two of \citet{eckhardt1985}'s consequences deserve to be stated here, because
this paper's headline results are empirical instances of both. First, they show
(\S V, Fig.~5) that \emph{certain intensity distributions make an $N$-version
system, on average, more prone to failure than a single component} --- redundancy
reducing reliability. \S\ref{sec:measure} reports a measured panel that does
exactly this. Second, they show that ``any error region for which a majority of
software components fail will limit the effectiveness of redundancy regardless of
the magnitude of the average component failure probability'', with the limiting
failure probability governed only by intensity mass at $\theta \geq 0.5$
(their Eq.~15). \S\ref{sec:quorum}'s irreducible floor $\lambda$ is a measured
instance of that region.

Our contribution here is therefore not the model, and not the phenomena --- both
were derived in 1985. It is that this model, proposed for human-written program
versions and rarely tested since, is the one that fits an automated verification
panel; that the alternative used in recent multi-agent and ensemble work does
not; and that the two predicted pathologies are observable on a real panel of
programs. We use ``difficulty'' where Eckhardt and Lee write ``intensity''; the
quantity is the same.

\paragraph{Aggregation and the independence discount.}
The Condorcet jury theorem under correlated votes has a substantial literature
\citep{ladha1992,boland1989}; the discount applied in practice is Kish's design
effect \citep{kish1965}, imported from cluster sampling and used on a judge panel
by \citet{kohli2026}. Beta-binomial
models for over-dispersed binary data are standard \citep{williams1975}; agreement
statistics such as Fleiss' $\kappa$ \citep{fleiss1971} summarise rater
concordance but do not, by themselves, answer how much independent evidence a
panel holds. Ensemble learning inherited this model rather than reinventing it, and the chain
matters because it ends at the diagnostic we use. \citet{hansen1990} carried
\citet{eckhardt1985} into classifier ensembles explicitly --- ``Eckhardt and Lee
have proposed a model incorporating such input specific difficulty within a study
of $N$-version programming'' --- and named the distribution over $\theta$ the
\emph{classification difficulty distribution}, which is Eckhardt and Lee's
intensity distribution $G$ under another name. \citet{kuncheva2003} then reduce
that distribution to a scalar diversity index, its variance, alongside nine
others. Closest of those to our setting is \citet{partridge1997}'s
\emph{coincident-failure diversity}, which is defined on software versions and
built directly on the failure-count distribution --- the same histogram as our
Figure~\ref{fig:kerr}. Their $\mathrm{CFD} = \sum_{n=1}^{N}\frac{N-n}{N-1}f_n$
weights a failure by how \emph{few} versions share it: it is $1$ when every
failure is unique to one version and $0$ when all $N$ fail together. Our
\S\ref{sec:quorum} floor is the mass they assign zero.

The distribution-of-errors diagnostic of \S\ref{sec:shape} is therefore
well-established. What this literature does with it is score a set on a scalar;
what we do is fit competing generative models to it and ask which reproduces a
measured panel. \citet{kuncheva2003} also report a finding we
echo: such measures track ensemble gain less reliably than intuition suggests ---
in their controlled simulations all ten relate similarly to majority-vote
improvement, but they close with doubts about the measures' usefulness for
building ensembles on real problems.

\paragraph{Review panels in practice.}
In three large open-source projects, review \emph{coverage} and
\emph{participation} were found to relate to post-release defect counts
\citep{mcintosh2014}; qualitative study of industrial review finds its benefits
extend well beyond defect-finding, to knowledge transfer and team awareness
\citep{bacchelli2013}. Neither measures reviewer \emph{head count} --- McIntosh
et al.'s participation metrics are self-approval, review speed, and discussion.
Head count is the cruder proxy practice has adopted, and what it is worth is the
subject of this paper. Automated judging with language models
\citep{zheng2023} made evaluation cheap enough to run at scale, and
execution-based benchmarks \citep{chen2021} supply exactly the kind of decidable
ground truth this study needs. What this paper adds to
\S\ref{sec:related-panels}'s contemporaneous measurement is a panel whose ground
truth is decided by execution, and a test of which dependence \emph{model}
reproduces it.

\section{The instrument}
\label{sec:instrument}

\subsection{A negative result first: LLM panels we could not measure}
\label{sec:e016}

Our first attempt deployed live LLM verifiers, because that is the panel
practitioners are actually building. Twenty verifiers --- four open-weight models
(\code{qwen2.5-1.5b}, \code{llama3.2-1b}, \code{gemma2-2b}, \code{smollm2-1.7b})
crossed with five prompt templates --- voted on all 72 candidates, 1440 votes
total.

The experiment failed, informatively. Of the 20 verifiers, six emitted a single
constant verdict across all 72 tasks and three more were $\geq 95\%$ one verdict.
Accuracy hid this: on a corpus that is $36.1\%$ viable, ``always reject'' scores
$0.639$. Youden's $J = \text{sensitivity} + \text{specificity} - 1$
\citep{youden1950} does not hide it. \citet{youden1950} states the relevant
property directly: $J$ ``has the value zero whenever a diagnostic test gives the
same proportion of positives for both diseased and control groups regardless of
what that proportion is'' --- so a rule that ignores the artifact scores zero
whichever way it is biased. Mean $J$ was $+0.0487$, and \textbf{zero of 20
verifiers had $J$ significantly above zero after Bonferroni correction}. The
20-member majority scored accuracy $0.514$ against $0.639$ for a constant
rejector, with a permutation test on panel $J$ giving $p = 0.13$.

What the votes \emph{did} respond to was the prompt: mean accept rate moved from
$0.251$ to $0.731$ across templates and $0.137$ to $0.799$ across models, while
carrying no task-level signal. The panel measured its own prompts.

We report this because it constrains the paper's scope and because it is a
methodological warning: \textbf{a correlation estimated over verifiers that do not
discriminate is a correlation between constants.} Any panel study must pass a
discrimination screen before its dependence structure means anything, and
accuracy on an unbalanced corpus is not that screen.

\subsection{A panel of partial oracles}

We therefore built verifiers that demonstrably verify. A verifier is a
\emph{partial test oracle}: it draws inputs from one named \emph{region} of a
problem's input domain --- \code{tiny}, \code{small}, \code{large},
\code{extreme}, \code{duplicate} --- and accepts a candidate only if it matches a
reference implementation on all drawn inputs. Five regions $\times$ five random
seeds gives $n = 25$ verifiers.

The corpus is 24 problems $\times$ 3 variants $=$ 72 candidate solutions. Ground
truth is not a judgement: each candidate is executed against hidden tests in a
subprocess and \code{viable} is the exit status. 26 are viable, 46 are not, a
base rate of $0.639$ non-viable. Labels reproduce exactly on regeneration
($72/72$).

\paragraph{Are the labels independent of the verifiers?}
They come from different artifacts. Ground truth is decided by executing
\emph{hand-written hidden tests} --- literal input/expected pairs such as
\code{("solve([1,2,3,4])", 2.5)} --- authored per problem from the specification.
The verifiers never see those tests: a partial oracle compares the candidate's
output against the executed output of the problem's \emph{reference
implementation} on inputs it draws itself. So the two sides are independent
encodings of one specification, and neither is derived from the other; expected
outputs are never hand-written for the oracles, and the hidden tests are never
generated from the reference. Two further guards matter. The variant labels
(``ok'', ``off\_by\_one'') are \emph{not} used as ground truth --- a mutant that
happens to pass every hidden test is recorded viable --- and inputs on which the
reference itself raises are treated as out of contract rather than as defects.
The residual shared dependency is the specification itself: a misconception
present in both the reference and the hand-written expectations would corrupt
both, and this design cannot detect that.

This panel has three properties that make it a usable instrument. Every verifier
is a program, so every error is a genuine missed defect rather than a
disagreement about taste. Every error vector is observed rather than assumed. And
region membership is a \emph{declared independence label} of exactly the kind
panel designers rely on --- ``these reviewers look at different things'' --- which
makes the label itself testable. \S\ref{sec:measure} tests it and finds it
wanting.

The panel passes the screen \S\ref{sec:e016} failed: \textbf{25 of 25 verifiers
discriminate above chance}, every Youden's $J$ significantly positive after
Bonferroni correction, at mean accuracy $p = 0.7956$.

\subsection{Effective independent panel size}
\label{sec:neff}

For independent verifiers, the probability that a panel's decision is wrong is a
binomial tail. With $\mathit{need} = \lfloor qn \rfloor + 1$ votes required to
accept at quorum $q$,
\begin{equation}
E_{\mathrm{indep}}(n, p, q) \;=\; \sum_{k=0}^{\mathit{need}-1} \binom{n}{k}\, p^{k}\,(1-p)^{n-k},
\end{equation}
which is strictly decreasing in $n$. We define the \textbf{effective independent
panel size} of a measured correlated panel as the size an \emph{independent}
panel would need to reproduce the same measured error rate,
\begin{equation}
\neff(n, p, \rho, q) \;=\; E_{\mathrm{indep}}^{-1}\!\left( \text{measured error rate} \right),
\end{equation}
interpolated linearly between bracketing odd panel sizes. The ratio $\neff/n$ is
the panel's \textbf{evidence efficiency}: the fraction of nominal reviewers that
survives as independent evidence.

\paragraph{The metric saturates at both ends, and this matters.}
$E_{\mathrm{indep}}$ is evaluated over odd panel sizes $1, 3, \dots, 201$, so
$\neff$ is \emph{clamped}: it returns $1$ for any measured error at or above the
single-verifier error, and $201$ for any error at or below the largest panel's.
The lower clamp is not cosmetic. A panel that performs \emph{worse} than one of
its members is not thereby ``worth one member'' --- it is worth less, and the
metric cannot say how much less. We therefore report $\neff$ only where it is
unclamped, and elsewhere report the error rates directly. Of the 315 grid cells
at $q=0.5$, 28 ($8.9\%$) sit on the lower clamp.

The metric is anchored where it is unclamped. On a simulated independent 5-panel
at $p = 0.75, q = 0.5$ the analytic error is $0.103516$ and the simulator
measures $0.103521$. Across the grid of \S\ref{sec:grid} the $\rho = 0$ column
recovers nominal panel size over a sevenfold range ($2.98$ to $20.99$) without
having been fitted (Figure~\ref{fig:neff}). The $\rho = 1$ column reads $1.00$
at every panel size, but that is the lower clamp firing rather than a
measurement, and we do not claim it as validation.

\section{How much independent evidence does the panel hold?}
\label{sec:measure}

\subsection{The declared label is real, and not nearly enough}

Pairwise error correlations, over all $\binom{25}{2} = 300$ pairs:

\begin{center}
\begin{tabular}{lrr}
\toprule
pair type & pairs & mean $\rho$ \\
\midrule
same region (declared dependent)   & 50  & $+0.8924$ \\
different region (declared independent) & 250 & $+0.5263$ \\
\midrule
all pairs & 300 & $+0.5873$ \\
\bottomrule
\end{tabular}
\end{center}

The declared label carries genuine signal: $+0.892$ within a region against
$+0.526$ across. But \textbf{verifiers that share no declared attribute still
share $53\%$ of their errors}. Treating a metadata boundary as an independence
boundary would overstate the panel's independent evidence by a wide margin. The
practical form of this is a measurement instruction, not a policy: before
weighting votes by declared groups, check whether cross-group error correlation
is actually lower than within-group correlation, and by how much.

\subsection{Headline: the panel is worse than its own members}

\begin{center}
\begin{tabular}{lrl}
\toprule
& error on 72 items & 95\% CI \\
\midrule
25-verifier majority              & $0.2083$ \;(15/72) & $[0.131, 0.316]$ \\
mean member                       & $0.2044$           & --- \\
\textbf{best member} (\code{large-s0}) & $\mathbf{0.1389}$ \;(10/72) & $[0.077, 0.237]$ \\
\midrule
\multicolumn{3}{l}{panel vs.\ best member: exact McNemar $p = 0.18$ \;($b=2$, $c=7$)} \\
\multicolumn{3}{l}{effective size $\neff$: \emph{below} the metric's floor of 1 (clamped)} \\
\multicolumn{3}{l}{design-effect heuristic $N/(1+(N-1)\rho)$: $1.66$} \\
\bottomrule
\end{tabular}
\end{center}

Under majority vote this panel does not merely fail to beat its members --- it
loses to them. \citet{eckhardt1985} predicted that this is possible for
sufficiently coincident intensity distributions (their Fig.~5, ``an intensity
distribution for which redundancy reduces reliability''); what follows is a
measurement of it on a panel of real programs. Aggregating 25 verifiers produces \emph{more} error than the mean
member and half again the error of the best one. Because the panel error exceeds
the single-verifier error, $\neff$ sits on its lower clamp
(\S\ref{sec:neff}): the honest statement is that the panel's effective
independent size is below $1$, not equal to it.

Two cautions belong with this number rather than in a later section. First, 72
items is a small sample: the panel--best-member gap is directional but not
significant ($p = 0.18$), and we claim the \emph{ordering}, not the magnitude.
Second, the design-effect heuristic reports $1.66$ where the measured value is
below $1$, so it overstates independence \emph{here}; \S\ref{sec:grid} shows when
that generalises and \S\ref{sec:related-panels} shows a measured panel where it
does not.

\subsection{When the heuristic is optimistic, and under what condition}
\label{sec:grid}

To separate ``this panel'' from ``panels'', we ran a full factorial simulation:
$n \in \{3,5,7,9,11,15,21\}$, $p \in \{0.60,0.70,0.75,0.80,0.90\}$,
$\rho \in \{0, 0.125, \dots, 1.0\}$, $q \in \{0.5, 0.7\}$, 100 seeds per cell ---
630 cells, $63{,}000$ runs. Figure~\ref{fig:neff} plots the resulting $\neff$
surface at $p = 0.75, q = 0.5$.

\begin{figure}[t]
\centering
\begin{tikzpicture}
\begin{axis}[
  width=0.62\linewidth, height=6.0cm,
  xlabel={verifier error correlation $\rho$},
  ylabel={effective panel size $\neff$},
  xmin=0, xmax=1, ymin=0, ymax=22,
  legend style={font=\scriptsize, at={(1.03,1)}, anchor=north west, draw=none},
  grid=major, grid style={gray!18},
  tick label style={font=\scriptsize}, label style={font=\small},
]
\addplot[thick, cReal, mark=*, mark size=1.2] table[x=rho,y=n21] {figures/neff-grid.dat};
\addlegendentry{$n=21$}
\addplot[thick, cBeta, mark=square*, mark size=1.2] table[x=rho,y=n15] {figures/neff-grid.dat};
\addlegendentry{$n=15$}
\addplot[thick, cShock, mark=triangle*, mark size=1.4] table[x=rho,y=n11] {figures/neff-grid.dat};
\addlegendentry{$n=11$}
\addplot[thick, cOne, mark=diamond*, mark size=1.4] table[x=rho,y=n7] {figures/neff-grid.dat};
\addlegendentry{$n=7$}
\addplot[thick, gray, mark=o, mark size=1.2] table[x=rho,y=n3] {figures/neff-grid.dat};
\addlegendentry{$n=3$}
\end{axis}
\end{tikzpicture}
\caption{Effective independent panel size against error correlation, at
$p=0.75$, $q=0.5$, from $63{,}000$ simulation runs. At $\rho = 0$ the metric
recovers nominal size across a sevenfold range ($2.98$ to $20.99$) without having
been fitted. The $\rho = 1$ column reads $1.00$ everywhere, but that is the
metric's lower clamp (\S\ref{sec:neff}), not a measurement. The curves converge
from $\rho \approx 0.5$ upward, where the spread across $n = 3 \dots 21$ falls
below $2.0$ effective verifiers; below that panel size still matters
substantially (spread $7.9$ at $\rho = 0.125$, $18.0$ at $\rho = 0$).}
\label{fig:neff}
\end{figure}

Two readings of this grid matter. The first is a ceiling: even $12.5\%$ error
correlation caps a 21-member panel at $10.6$ effective verifiers, and at
$\rho = 0.5, p = 0.75$ going from 15 to 21 reviewers buys $+0.05$ effective
verifiers. The second concerns the heuristic. Re-deriving each cell under the two
competing dependence models of \S\ref{sec:shape} gives:

\begin{center}
\begin{tabular}{lr}
\toprule
cells where $N/(1+(N-1)\rho)$ overstates independence & \\
\midrule
under the shared-shock model & $17/441$ \; ($4\%$) \\
under the per-item difficulty model & $\mathbf{441/441}$ \; ($100\%$) \\
\bottomrule
\end{tabular}
\end{center}

Under the item-difficulty shape the design-effect heuristic overstates available
independence in \emph{every} cell. The rule this licenses is a conditional one,
and the condition matters: \textbf{under item-difficulty dependence,
$N/(1+(N-1)\rho)$ overstates a panel's independent evidence --- so establish the
dependence shape before using it to size a panel.} We deliberately do not state
this unconditionally. The $441/441$ figure is computed \emph{under} the
item-difficulty model; comparing it against a heuristic derived under different
assumptions demonstrates that two models disagree, not that the heuristic is
false in general. \S\ref{sec:related-panels} describes a measured panel where it
is in fact accurate.

This also corrects a conclusion we had drawn ourselves. Under shared shock,
$\neff$ saturates as the panel grows and does so most sharply for accurate
verifiers, suggesting that high accuracy is a reason to stop adding reviewers.
Under per-item difficulty the regimes invert. Extending past the grid at
$p = 0.90, \rho = 0.125$:

\begin{center}
\begin{tabular}{lrrrrrr}
\toprule
$n$ & 5 & 11 & 21 & 41 & 81 & 151 \\
\midrule
shared shock        & $3.82$ & $4.57$ & $4.59$ & $4.59$ & $4.59$ & $4.59$ \\
per-item difficulty & $2.82$ & $3.93$ & $4.56$ & $4.89$ & $5.16$ & $5.40$ \\
\bottomrule
\end{tabular}
\end{center}

Shared shock pins at $4.59$ and never moves; per-item difficulty crosses it near
$n = 21$ and keeps climbing. Conversely at $p = 0.75, \rho = 0.5$ the per-item
model flattens almost immediately, at $\neff \approx 1.6$ rather than the $4.1$
shared shock reports. \textbf{High correlation is what caps a panel, and it caps
it harder and sooner than the shared-shock model says; high accuracy is not
itself a reason to stop adding verifiers.}

\subsection{A contemporaneous panel where the heuristic holds}
\label{sec:related-panels}

While this work was in preparation, \citet{kohli2026} reported the same headline
on a very different panel: nine frontier LLM judges from seven model families, on
three natural-language-inference benchmarks with 100 human annotations per item.
That panel is worth $\neff = 2.18$ of 9 (bootstrap CI $[2.07, 2.31]$), the best
single judge matches or beats the full panel, and the conclusion --- that
correlated judges, not the aggregation rule, are the binding constraint --- is
ours as well. Agreement between two independent instruments, one on code with
executable ground truth and one on natural language with human labels, is the
strongest evidence either paper offers that the effect is real rather than an
artifact of one construction. That work has priority on the effective-votes
framing for judge panels, and we make no first-measurement claim.

The two panels disagree on exactly two points, and the disagreement is
informative.

\textbf{The heuristic.} \citet{kohli2026} finds $N/(1+(N-1)\rho)$ \emph{accurate}
on their panel: the empirical $\neff(k)$ curve tracks $k/(1+(k-1)\cdot 0.391)$,
and an independent eigenvalue estimate ($2.16$) matches the Kish value ($2.18$).
We find it optimistic on ours. The panels differ in measured correlation
($0.391$ against $0.5873$) and, more consequentially, in dependence shape.

\textbf{The shape.} Their permutation test rejects shared item difficulty as the
explanation of their inter-judge correlation ($p < 10^{-4}$; only $6.8$--$13.5\%$
of the gap is attributable to it). We find item difficulty is the shape that fits
ours. The honest reading favours their result as a constraint on ours rather than
the reverse, because \emph{our panel's item-difficulty structure is
manufactured}: five input regions $\times$ five seeds induces per-item difficulty
almost by construction, since a candidate that fails on \code{extreme} inputs
fails for every verifier drawing from that region. Their panel's diversity is not
designed at all --- it is whatever seven model families happen to share.

The defensible summary is therefore not that one paper is right, but that
\textbf{dependence shape is a property of how a panel was assembled and must be
measured per panel rather than assumed}. That is a weaker claim than either paper
makes alone, and better supported than either.

\section{The shape of verifier dependence}
\label{sec:shape}

\subsection{A single $\rho$ does not reproduce the panel}

Feeding the panel's own measured correlation back into the shared-shock mixture
--- with probability $\rho$ all verifiers share one correctness state, otherwise
they are independent --- under-predicts what the panel does:

\begin{center}
\begin{tabular}{lrr}
\toprule
model & panel error & vs.\ real \\
\midrule
real 25-verifier majority        & $0.2083$ & --- \\
independent ($\rho = 0$)         & $0.0005$ & $417\times$ too low \\
flat shared shock at $\rho=0.587$ & $0.1216$ & $1.71\times$ too low \\
nested shared shock ($g=0.526$, $b=0.773$) & $0.1290$ & $1.62\times$ too low \\
\bottomrule
\end{tabular}
\end{center}

A nested variant matching \emph{both} the within-region and cross-region
correlations barely improves on the flat one, and heterogeneous accuracy does not
close the gap either: re-simulating with each verifier's own measured accuracy
gives $0.1228$, still $0.086$ short. The problem is not the parameter value.

\subsection{Real panels fail partially; shared shock cannot}

The diagnosis is visible in the distribution of how many verifiers err per task
(Figure~\ref{fig:kerr}). That diagnostic is not new --- it is the empirical form
of \citet{eckhardt1985}'s intensity distribution, called the classification
difficulty distribution by \citet{hansen1990} and reduced to a scalar index by
\citet{kuncheva2003}. What has not been done, and is what
follows, is to treat the distribution as something to be \emph{explained}: to fit
competing generative models to it at matched parameter count and ask which one
reproduces a real panel.

\begin{figure}[t]
\centering
\begin{tikzpicture}
\begin{axis}[
  name=full,
  width=0.92\linewidth, height=4.4cm,
  xlabel={}, ylabel={tasks (of 72)},
  xmin=-0.7, xmax=25.7, ymin=0, ymax=46,
  xtick={0,5,10,15,20,25},
  ybar=0pt, bar width=2.4pt,
  legend style={font=\scriptsize, at={(0.5,1.02)}, anchor=south,
                legend columns=3, draw=none, column sep=6pt},
  grid=major, grid style={gray!18},
  tick label style={font=\scriptsize}, label style={font=\small},
]
\addplot[fill=cReal, draw=cReal] table[x=k,y=observed] {figures/kerr-distribution.dat};
\addlegendentry{observed}
\addplot[fill=cBeta!70, draw=cBeta] table[x=k,y=betabinom] {figures/kerr-distribution.dat};
\addlegendentry{per-item difficulty}
\addplot[fill=cShock!60, draw=cShock] table[x=k,y=sharedshock] {figures/kerr-distribution.dat};
\addlegendentry{shared shock}
\node[font=\scriptsize\itshape, anchor=north east] at (axis cs:25.4,44) {full scale};
\end{axis}

\begin{axis}[
  at={(full.below south west)}, anchor=north west, yshift=-0.9cm,
  width=0.92\linewidth, height=4.0cm,
  xlabel={number of the 25 verifiers wrong on a task},
  ylabel={tasks (of 72)},
  xmin=12.3, xmax=24.7, ymin=0, ymax=6.2,
  xtick={13,15,17,19,21,23,24},
  ybar=0pt, bar width=3.4pt,
  grid=major, grid style={gray!18},
  tick label style={font=\scriptsize}, label style={font=\small},
]
\addplot[fill=cReal, draw=cReal] table[x=k,y=observed] {figures/kerr-distribution.dat};
\addplot[fill=cBeta!70, draw=cBeta] table[x=k,y=betabinom] {figures/kerr-distribution.dat};
\addplot[fill=cShock!60, draw=cShock] table[x=k,y=sharedshock] {figures/kerr-distribution.dat};
\node[font=\scriptsize\itshape, anchor=north east] at (axis cs:24.5,5.9)
      {partial panel failures ($13 \le k \le 24$), magnified};
\end{axis}
\end{tikzpicture}
\caption{Where each model puts its mass, at full scale (top) and magnified over
the partial-failure region (bottom). A 25-verifier majority fails when
$k \ge 13$; it fails \emph{partially} when $13 \le k \le 24$, a majority wrong
while a minority is still right. That is how 11 of the panel's 15 real failures
happen, and shared shock is flat against zero across the whole region
($0.01$ tasks expected) --- conditional on the shock it produces unanimity,
conditional on no shock a tight binomial spike near $k = 5$ that dies well
before 13. It compensates by over-predicting unanimous failure ($8.5$ against
$4$ observed). The per-item difficulty model reproduces the region at the same
parameter count ($11.1$ predicted against $11$ observed).}
\label{fig:kerr}
\end{figure}

Replacing ``with probability $\rho$ everyone shares one correctness state'' with
``each task has a difficulty $d \sim \mathrm{Beta}(\alpha,\beta)$ and each
verifier errs independently with probability $d$'' gives a beta-binomial
\citep{williams1975}. This is \citet{eckhardt1985}'s model with a Beta mixing
distribution: their Corollary writes system failure as
$p_N = \int h(y;N)\,dG(y)$, a binomial mixed over the \emph{intensity
distribution} $G$, and they deliberately leave $G$ unspecified. Choosing
$G = \mathrm{Beta}$ is our parametric commitment, not theirs --- a point
\S\ref{sec:threats} already concedes. It has the same two free parameters as flat shared
shock --- a mean and a correlation --- so the comparison is of \emph{shape}, not
of freedom:

\begin{center}
\begin{tabular}{lrrr}
\toprule
 & real & per-item difficulty & shared shock \\
\midrule
panel error at majority & $0.2083$ & $0.1847$ & $0.1216$ \\
partial failures ($13 \le k \le 24$) & $11$ & $11.1$ & $\mathbf{0.01}$ \\
unanimous failures ($k = 25$) & $4$ & $2.3$ & $8.5$ \\
implied pairwise $\rho$ & $0.5873$ & $0.5713$ & --- \\
\bottomrule
\end{tabular}
\end{center}

The fitted correlation lands on the independently measured $\rho$, the partial
failure count is reproduced almost exactly, and the panel error is $3.7\times$
closer. Generalising over the $441$-cell grid, the per-item model predicts more
error than shared shock in $438/441$ cells, by a median factor of $1.27\times$ and
up to $2.71\times$. \textbf{Recommendation: parameterise verifier dependence by an
item-difficulty distribution, not by a shared-shock probability. $\rho$ remains a
useful summary of a panel; it is not a sufficient statistic for its error.}

\section{The quorum frontier and the irreducible blind spot}
\label{sec:quorum}

Majority vote is one point on a curve. Sweeping the acceptance threshold from 1
to 25 votes gives the panel's full frontier (Figure~\ref{fig:quorum}).

Before reading anything off that curve, it must clear the screen
\S\ref{sec:e016} demanded of the LLM panel: a verification panel must beat a rule
that ignores the code. On this corpus ``reject everything unread'' errs on the 26
viable candidates, a rate of $0.3611$. \textbf{The panel does not beat that
baseline until $\mathit{need} = 6$}, and at $\mathit{need} = 5$ it exactly ties it.
Majority ($0.2083$) and the whole high-quorum region clear it comfortably, but the
left third of the frontier does not, and nothing should be concluded from that
region. We apply the test to ourselves because \S\ref{sec:e016} applies it to
others.

\begin{figure}[t]
\centering
\begin{tikzpicture}
\begin{axis}[
  width=0.92\linewidth, height=6.2cm,
  xlabel={votes required to accept ($\mathit{need}$, of 25)},
  ylabel={panel error rate},
  xmin=1, xmax=25, ymin=0, ymax=0.58,
  legend style={font=\scriptsize, at={(0.98,0.97)}, anchor=north east, draw=none},
  grid=major, grid style={gray!18},
  tick label style={font=\scriptsize}, label style={font=\small},
]
\addplot[very thick, cReal] table[x=need,y=real] {figures/quorum-frontier.dat};
\addlegendentry{measured panel}
\addplot[thick, cShock, dashed] table[x=need,y=sharedshock] {figures/quorum-frontier.dat};
\addlegendentry{shared shock (2 params)}
\addplot[thick, cBeta, dotted] table[x=need,y=betabinom] {figures/quorum-frontier.dat};
\addlegendentry{per-item difficulty (2 params)}
\addplot[thick, cOne, dashdotted] table[x=need,y=oneinflated] {figures/quorum-frontier.dat};
\addlegendentry{one-inflated (3 params)}
\addplot[black!60, thick, dashed, domain=1:25, samples=2] {0.3611};
\addlegendentry{``reject everything unread''}
\addplot[gray, thin, domain=1:25, samples=2, forget plot] {0.0556};
\node[font=\scriptsize, gray, anchor=south west] at (axis cs:1.4,0.062) {blind-spot floor $\lambda = 0.0556$};
\end{axis}
\end{tikzpicture}
\caption{The real quorum frontier against three models, with the trivial
baseline. RMSE over all 25 thresholds: shared shock $0.0766$, per-item difficulty
$0.0251$, one-inflated $0.0328$. From $\mathit{need} = 17$ upward shared shock is
frozen at $0.1186$ while reality keeps falling; past $\mathit{need} = 22$ reality
stops at $0.0556$ while the per-item model keeps promising improvement. Note the
left of the plot: the panel does not beat ``reject everything unread'' until
$\mathit{need} = 6$, and nothing should be concluded from that region.}
\label{fig:quorum}
\end{figure}

At unanimity the three numbers separate cleanly, and neither two-parameter model
lands on the truth:
\[
\text{real} = 0.0556, \qquad
\text{shared shock} = 0.1186 \;(2.13\times\text{ high}), \qquad
\text{per-item} = 0.0313 \;(1.77\times\text{ low}).
\]

The disagreement is structural rather than a fitting artifact.

\begin{itemize}
\item \textbf{Shared shock} puts $P(\text{all } n \text{ wrong}) = \rho\mu + (1-\rho)\mu^{n}$,
  tending to a hard floor $\rho\mu = 0.1186$ at \emph{every} panel size. The model
  asserts that past some threshold, raising the quorum cannot help at all.
\item \textbf{Per-item difficulty} puts $P(\text{all } n \text{ wrong}) = \mathbb{E}[d^{n}] = B(\alpha+n,\beta)/B(\alpha,\beta)$,
  which decays as $n^{-\beta}$ with $\beta = 0.5762$ --- slowly, but to zero. The
  model asserts there is no floor, and that enough verifiers will eventually
  verify anything.
\item \textbf{Reality has a floor, but not shared shock's.} It sits at
  $\lambda = 0.0556$: four defects that \emph{every one} of the 25 verifiers
  misses. That is a shared blind spot of the panel, not a correlated shock. It is
  also the region \citet{eckhardt1985} identified as the binding constraint ---
  intensity mass at $\theta \geq 0.5$, which caps achievable reliability
  ``regardless of the magnitude of the average component failure probability''.
\end{itemize}

The distinction is operationally important. A shared shock is a property of the
\emph{correlation}; a blind spot is a property of \emph{which verifiers you
chose}. You escape the first by decorrelating an existing panel, and the second
only by adding a verifier of a different kind. Adding one-inflation --- a
$\lambda$-atom on ``all wrong'' plus a beta-binomial over the rest --- captures
both, at the cost of a third parameter. Once the blind spot is represented
explicitly, the correlation remaining in the reducible units is $0.4513$ rather
than the marginal $0.5873$; feeding the marginal value \emph{and} arming the atom
would double-count the same shared failures.

Finally, the assumed shape --- not the correlation --- decides the aggregation
rule. Sweeping the cost-optimal threshold $q^{*}$ over correlation, base rate and
false-accept/false-reject cost ratio, at $n = 25$ the optimum spans $12$--$14$
under shared shock but $1$--$25$ under per-item difficulty, with shifts as large
as $\pm 11$ votes in the same cell. Under shared shock, tuning the quorum on this
panel appears to buy nothing ($1.00\times$); on the real panel it is the single
largest available error reduction ($3.75\times$, from $0.2083$ at majority to
$0.0556$).

\section{Consequences for panel design}
\label{sec:consequences}

\subsection{Group balancing is not a safe default}

If declared groups are the problem, weighting by them looks like the fix: give
each declared independence group one vote instead of counting every verifier
equally. Under shared shock this crosses over from harmful to helpful at
$\rho = 0.25$ for a reference $[7,1,1,1,1]$ panel at $p = 0.75$, and that
crossover survives the change of shape --- it sits at $\rho = 0.25$ under both
models, a robustness result.

It does not survive the groups being fake. When task difficulty is shared across
groups rather than within them --- the regime the measured panel's
$0.526$ cross-group correlation points toward --- group balancing \textbf{never
wins at any correlation}, and loses by up to 3 percentage points:

\begin{center}
\begin{tabular}{lrrl}
\toprule
$\rho$ & naive majority & group-balanced & winner \\
\midrule
$0.000$ & $\mathbf{0.0345}$ & $0.0659$ & naive \\
$0.125$ & $\mathbf{0.1200}$ & $0.1366$ & naive \\
$0.250$ & $\mathbf{0.1719}$ & $0.1809$ & naive \\
$0.500$ & $\mathbf{0.2232}$ & $0.2255$ & naive \\
$0.750$ & $\mathbf{0.2448}$ & $0.2449$ & naive \\
$1.000$ & $0.2492$ & $0.2492$ & tie \\
\bottomrule
\end{tabular}
\end{center}

The mechanism is straightforward: balancing pays a real cost, discarding the
extra votes inside the large group, and that cost is only worth paying when the
small groups supply evidence the large group lacks. When every group errs on the
same hard items, the discarded votes are pure loss. At $\rho = 1$ both rules
report $\approx 1 - p$: the panel has collapsed to a single verifier and no
aggregation rule can recover anything.

\subsection{Correlation, not accuracy, is the exploitable surface}
\label{sec:adversarial}

A last consequence concerns adversarial pressure, and we state its status
plainly: this is a simulation result built on the measured panel's parameters,
not a second measurement. In a matched-budget search benchmark, contributors
optimise to pass the gate --- drawing $k$ candidates and submitting the one that
looks best on exactly the quantity the gate ranks. Holding per-verifier accuracy
fixed at the measured $p = 0.7956$ and the head count fixed at 25, we vary only
the \emph{dependence structure}. Because a panel's dependence has two components
that our own \S\ref{sec:quorum} insists are different mechanisms, we run each
alone rather than moving both together:

\begin{center}
\begin{tabular}{lrrr}
\toprule
panel ($n=25$, $p=0.7956$) & correlation & blind spot & catastrophic seeds \\
\midrule
\code{independent}      & $0.0$    & $0.0$    & $0$ \\
\code{blindspot\_only}  & $0.0$    & $0.0556$ & $14$ \\
\code{correlated\_only} & $0.4513$ & $0.0$    & $\mathbf{53}$ \\
\code{measured}         & $0.4513$ & $0.0556$ & $58$ \\
\bottomrule
\end{tabular}
\end{center}

Catastrophic seeds are out of 100, for the best-performing search strategy, at
the highest adversary fraction ($0.2$) and effort ($k=8$). The two end rows
reproduce the published matrix exactly; the two middle rows are the
single-factor cells this paper adds.

\textbf{Correlation carries the effect.} Correlation alone reproduces $53$ of the
$58$ catastrophic seeds --- $91\%$ --- while the blind spot alone reproduces
$14$. The same head count and the same per-verifier accuracy separate immunity
from near-coin-flip failure, and it is the correlation component that does it.
Two honest caveats: the components are \emph{sub-additive} ($53 + 14 = 67$
against a joint $58$), so they overlap rather than compose, and a blind spot of
$0.0556$ is a much smaller perturbation than a correlation of $0.4513$, so this
compares the two effects \emph{as measured on our panel} rather than at matched
intensity. What the decomposition rules out is the alternative reading --- that
the gap is really the blind spot, with correlation along for the ride.

A second observation from the same benchmark is worth recording because it is a
measurement trap. A panel that errs at the same rate for everyone is not a panel
that harms everyone equally: an error only hurts in the direction a producer is
exposed to. Strategies whose proposals are almost always viable (base viability
$\geq 0.89$) can essentially only suffer false \emph{rejects}; a strategy
proposing viable work $40\%$ of the time is mainly exposed to false
\emph{accepts}. Across five strategies spanning a base-viability range of
$0.588$, the least viable one took nine times the utility damage of the next
worst. Its \emph{accept rate went up} as the gate degraded, from $0.398$ to
$0.488$, while every healthy strategy's fell. \textbf{Reading accept rate as gate
health inverts the ranking: the strategy whose throughput improved was the one
being destroyed.}

\section{Threats to validity}
\label{sec:threats}

\paragraph{The panel's diversity is constructed.} This is the primary threat. Our
25 verifiers are 5 input regions $\times$ 5 seeds, a diversity structure we
designed so that the declared independence label would be testable. It is not a
sample of the diversity structures found in deployed review panels, and the
measured $\rho = 0.5873$ should be read as one real value, not as an estimate of
a population parameter. The qualitative findings --- that a declared label
under-predicts dependence, that failures are partial, that a blind spot exists ---
are the transferable claims; the specific numbers are not.

\paragraph{The per-item difficulty model is validated against one panel.} It fits
far better than shared shock, which is a low bar. A Beta difficulty distribution
is a modelling choice --- the natural two-parameter conjugate form --- and not a
shape we measured. Both models also treat verifiers as exchangeable, with
identical accuracy and one common correlation; real panels are heterogeneous, and
while we showed heterogeneity does not by itself explain the shared-shock gap,
neither model represents it.

\paragraph{Corpus scale and domain.} 72 candidates from 24 problems, all
self-contained algorithmic programs with executable hidden tests. The blind-spot
floor rests on four candidates. The corpus supports a base rate ($0.639$
non-viable) that would differ in a repository where most submitted work is
roughly correct, and \S\ref{sec:quorum}'s cost-optimal thresholds are explicitly
base-rate dependent.

\paragraph{The LLM result is about small models.} \S\ref{sec:e016}'s panel used
four open-weight models in the 1--2B parameter range. It establishes that
verifier competence must be screened before dependence is estimated; it does not
establish that larger models fail the screen, and we expect they do not.

\paragraph{The adversarial results are simulated.}
\S\ref{sec:adversarial}'s decomposition is generated by a search benchmark whose
verification panel is parameterised from our measurements, not observed in a live
system. The single-factor cells separate correlation from blind spot
\emph{within that simulator}; they do not establish which component an attacker
on a real review gate would exploit. In that
simulator the panel reads each candidate through a single viability bit, so
questions of the form ``can an attacker learn what the panel cannot see'' are
unanswerable within the instrument: the blind spot has no address in artifact
space. We report the contrast as a consequence of the measured parameters, not as
evidence about real attackers.

\paragraph{Quorum restriction in the grid.} The $441$-cell comparison of
\S\ref{sec:grid} is restricted to $q = 0.5$, where ``the panel is wrong'' and
``the panel wrongly accepts'' coincide. Above $0.5$ they diverge, and separating
false accepts from false rejects requires the corpus base rate --- which is why
the quorum analysis of \S\ref{sec:quorum} is done directly on the measured panel
rather than on the grid.

\section{Reproduction}
\label{sec:repro}

All artifacts are committed in the IDKMesh repository \citep{idkmesh}. The panel
votes, the corpus with its execution-decided labels, and the simulation grid are
under \code{experiments/results/}; the analyses are pure standard-library Python.

\begin{quote}\small
\begin{verbatim}
python3 sim/e017_verify.py --seeds 5 --out e017-votes.jsonl   # regenerate votes
python3 sim/e017_analyze.py experiments/results/E017-partial-oracle-votes.jsonl.gz
python3 sim/e018_dependence_models.py    # the 441-cell model comparison
python3 sim/e020_quorum_frontier.py      # the quorum frontier, deterministic
python3 paper/make_figures.py            # regenerate every figure in this paper
PYTHONPATH=. python3 paper/deconfound_panel.py \
    --seeds 100 --agents 64 --generations 50 --change-at 25 --jobs 4
\end{verbatim}
\end{quote}

The figures are generated by importing the same analysis modules the experiment
records use, so no number in this paper can drift from the artifact it came from.
The last command runs \S\ref{sec:adversarial}'s single-factor cells; it registers
the two extra panels on the benchmark's own panel table at run time, so it adds
cells without modifying any committed experiment module.
Full experiment records, including the preregistered predictions that were
falsified, are in \code{experiments/E015}--\code{E020}.

\section{Conclusion}

A verification panel's head count is a budget, not a measurement. On a panel
where every member demonstrably works and every error is a genuine missed defect,
25 verifiers delivered less evidence than the best one of them, the standard
independence discount overstated what was available \emph{on this panel}, and the
dependence model in common use was wrong in shape rather than in parameter --- unable to produce the partial failures that
are how real panels actually fail. Beneath all of that sat a floor no aggregation
rule could reach: defects invisible to every verifier in the panel, escapable
only by adding a verifier of a different kind.

The operational summary is short. Screen your verifiers for discrimination before
you trust any statistic about them. Estimate independence from observed error
vectors, never from declared diversity. Establish the dependence \emph{shape}
before applying $N/(1+(N-1)\rho)$, which overstates independence when difficulty
is shared across items. Check your panel against a rule that ignores the
artifact. Measure your quorum, since it is likely a larger lever than your panel
size. And treat the dependence structure, not accuracy, as the property an
adversary will find. Every one of these is a measurement instruction, which is
the paper's real conclusion: a panel's evidentiary value is not a property you
can assume from its size, its members' accuracy, or its declared diversity.

\paragraph{Data availability.} Code, artifacts, and experiment records:
\url{https://github.com/MSKazemi/idkmesh}.

\paragraph{AI/tool provenance.} Drafting and revision of this manuscript, and the
analysis and figure-generation code in \S\ref{sec:repro}, were materially
assisted by an AI coding assistant (Claude Code). Every quantitative claim
reported here was independently re-derived from the committed artifacts by
re-running the cited scripts rather than accepted from prose, and the
manuscript was revised against a prior adversarial review
(\S\ref{sec:repro}'s repository carries that review alongside this source).
AI assistance is not represented as independent human or external scientific
review, and this manuscript has not undergone peer review at a venue as of
this revision.

\bibliographystyle{plainnat}
\bibliography{refs}

\end{document}
